% =====================================================================
% SEC 02 — Trabalhos Relacionados
% =====================================================================
\section{Trabalhos Relacionados}
\label{sec:rel}

\subsection{Ranqueamento e recuperação}

A recuperação \emph{lexical} clássica, exemplificada pelo BM25
\citep{RobertsonZaragoza2009BM25}, otimiza a relevância por sobreposição de
termos e é amplamente usada como \emph{baseline} de recuperação. Métodos densos
de \emph{passage retrieval} \citep[cf.][]{KhattabZaharia2020ColBERT} representam
trechos e consultas em espaço vetorial latente, capturando semelhança semântica;
contudo, ambos tendem a concentrar resultados em um subespaço homogêneo, o que
reduz a diversidade do conjunto recuperado \citep{He2023RethinkingRetrieval}.

\subsection{Retrieval-Augmented Generation}

O paradigma RAG \citep{Lewis2020RAG} acopla recuperação a geração ancorada.
Abordagens modernas exploram o enriquecimento de consultas e a avaliação
\emph{adaptativa} \citep{Jeong2024AdaptiveRAG}, a intercalação de passagens
\citep{Trivedi2023Interleaving}, a correção reflexiva \citep{Yan2024CRAG} e o
uso de \emph{grafos de conhecimento} \citep{Edge2024GraphRAG}. Na maioria desses
trabalhos, o foco está em melhorar a \emph{qualidade da evidência recuperada};
a \emph{diversidade} do conjunto raramente é tratada como primeiro-cidadão,
especialmente em termos de métrica explícita.

\subsection{Diversificação em Recuperação de Informação}

O \emph{Maximal Marginal Relevance} (MMR) \citep{CarbonellGoldstein1998MMR} é o
arcabouço clássico de diversificação: seleciona itens maximizando uma combinação
linear de relevância e dissimilaridade em relação aos já selecionados,
\[
\label{eq:mmr}
\mathrm{MMR}(d) \;=\; \lambda\,\mathrm{relev}(d) \;-\; (1-\lambda)\,
\max_{d_j\in S} \Sim(d, d_j) \; ,
\]
onde $\lambda\in[0,1]$ controla o balanço. O MMR é eficaz, porém \emph{sensível a
parâmetros} ($\lambda$) e à escolha da função de similaridade, e tem custo
$O(K\cdot|D|)$. A proposta pós-Recamán \citep{Alekseyev2022RecamanCousins} situa-se
na mesma família, mas substitui a otimização gulosa por \emph{offsets
determinísticos} derivados da sequência de Recamán, removendo parâmetros e
aleatoriedade. A pergunta que este artigo coloca é se essa escolha preserva a
capacidade de diversificar \emph{conteúdo} (âncora) em cenários realistas.

\subsection{Diversidade adaptativa e a limitação do $\lambda$ fixo}

Uma limitação recorrente do MMR é a \emph{escolha do $\lambda$}. Diferentes
estratégias foram propostas para superá-la. \emph{Probabilistic Latent MMR}
(PLMMR) \citep{GuoSanner2010PLMMR} deriva uma variante de MMR por inferência
variacional num modelo de grafos latentes que \emph{elimina a necessidade de
balançar manualmente} relevância e diversidade, mostrando ainda que a diversidade
de um conjunto deve ser \emph{relevante à consulta}. Na mesma direção,
\emph{Learning MMR} \citep{Xia2015LearningMMR} aprende o modelo de relevância
otimizando diretamente métricas de diversidade a partir de dados. Mais
recentemente, no contexto de RAG, \emph{DF-RAG} \citep{Khan2026DFRAG} mostrou
empiricamente (em cinco \emph{benchmarks} de QA) que \textbf{o $\lambda$ ótimo
varia por query e por dataset}, e que um \emph{Oracle} que escolhe o $\lambda$
ideal por query supera o RAG vanilla em até $\sim 18\%$ de F1 absoluto — dos quais
o próprio DF-RAG captura até $91{,}3\%$ adaptando o $\lambda$ \emph{em tempo de
teste}, sem fine-tuning, via um \emph{Planner/Evaluator} LLM. Já \emph{Vendi-RAG}
\citep{RezaeiDieng2025VendiRAG} adapta iterativamente um parâmetro de diversidade
$s$ com base na qualidade da resposta avaliada por um juiz LLM, usando a
\emph{Vendi Score} como métrica explícita de diversidade de conjuntos.

Nosso trabalho dialoga com essa linha \emph{sem depender de LLM}: o
$\lambda$ adaptativo do HABD é derivado por uma \emph{heurística determinística}
(índice de mistura de âncoras no topo do ranking), o que o torna auditável,
reprodutível e de custo desprezível. Em contraste com DF-RAG e Vendi-RAG, não
treinamos nem consultamos um modelo para escolher o balanço; medimos a
\emph{estrutura} do ranking (monopolista vs. misto) e ajustamos o balanço em
consequência. Isso reforça a tese comum — $\lambda$ fixo é subótimo — e oferece
uma alternativa determinística para cenários sem acesso a LLM de julgamento.

\subsection{Posição deste trabalho}

Ao contrário de trabalhos que propõem novos diversificadores e assumem ganho,
este artigo submete as propostas a um \emph{turno empírico controlado},
comparando-as ao \emph{top-k} e ao MMR em métricas explícitas. Reportamos um
resultado em que a proposta posicional \emph{empatia} com o \emph{top-k} e em que
o HABD, apesar de superar em diversidade/cobertura, \emph{não demonstra} ganho
dentro da tolerância de relevância — vereditos \texttt{refuta\_H2} e
\texttt{refuta\_H3}. Essa postura de \emph{reporte negativo/quilibrado}, ancorada
não apenas na teoria, mas na \emph{medição explícita} de diversidade
\citep{Yu2024RAGEvalSurvey}, é rara na literatura e contribui para a robustez da
área.
